% Chapter 4. Formulae are direct translations of original Core.lean.
\chapter[Supercharges and Source Core Form]{Sixteen Supercharges and the Normalized Smooth-Core Form}
\label{ch:charges}

\section{Why we begin with charges instead of guessing a Hamiltonian}
The supersymmetric structure of this model is encoded in sixteen first-order
operators, each carrying a spatial spinor index. Starting from these
operators is important for two reasons. First, the fermionic interaction,
color structure, and normalization arise from specific finite sums: replacing
one of them by a familiar-looking textbook expression can change a sign or
factor. Second, defining a positive energy form directly from the charge
column gives a controlled route to a closed physical Hamiltonian, avoiding
an unjustified choice among self-adjoint realizations of a formal second-order
differential expression.

We retain the source's \emph{Hermitian color basis} $t_A$, satisfying
\begin{equation}\label{eq:color-basis}
 t_A^\dagger=t_A,\qquad \operatorname{tr}(t_A)=0,
 \qquad \operatorname{tr}(t_A t_B)=\delta_{AB}.
\end{equation}
For two color matrices define the source bracket and structure constants by
\begin{equation}\label{eq:source-bracket}
 [X,Y]_{\mathrm{s}}:=-i(XY-YX),\qquad
 f_{ABC}:=\Re\operatorname{tr}\bigl(t_A[t_B,t_C]_{\mathrm{s}}\bigr).
\end{equation}
The source data also require that the $t_A$ span the real vector space of
Hermitian traceless color matrices. When specialized to the source's
normalized Pauli basis, these definitions fix the numerical factors in the
familiar $SU(2)$ structure constants. We do not reset that normalization
by convention in a downstream argument.

\section{The literal differential expressions}
For $i=0,\ldots,8$ and $A=0,1,2$, write
$\partial_{iA}:=\partial/\partial x_i^A$, and use the spinor matrices
$\gamma_i$ and $\gamma^{ij}$ of \eqref{eq:gamma-clifford} and
\eqref{eq:gamma-two}. For $\alpha\in\mathsf{Spin}$ define on
$\mathscr D=C_c^\infty(\mathbb R^{27};\mathcal F)$
\begin{align}\label{eq:source-charge}
 Q_\alpha \psi
 ={}&-i\sum_{i,A,\beta}
   (\gamma_i)_{\alpha\beta}\,\theta_{\beta A}\,
   \partial_{iA}\psi \\
 &+\frac12\sum_{i,j,A,B,C,\beta}
 f_{ABC}\,(\gamma^{ij})_{\alpha\beta}
 \theta_{\beta A}\,x_i^B x_j^C\,\psi .\nonumber
\end{align}
These finite sums are \emph{ordered}, as in the original source
\texttt{AlgebraData.charge}. The factor $1/2$ in the second term is
part of the definition. Replacing ordered pairs with
unordered pairs without tracking this coefficient would give an incorrect
source operator.

The matrices $\theta_{\beta A}$ are bounded on the \emph{finite} fiber
$\mathcal F$, but $\partial_{iA}$ and multiplication by $x_i^B x_j^C$
are unbounded after passage to $\mathscr H$. Thus \eqref{eq:source-charge}
initially defines a linear map $\mathscr D\to\mathscr D$ and, after
$\iota$, a densely defined map into the $L^2$ output. It is premature to
call its formal adjoint its Hilbert-space adjoint until the operator domain
has been constructed.

\begin{definition}[Original charge column and smooth-core form]
The original sixteen-charge column on $\mathscr D$ is
\begin{equation}\label{eq:charge-column}
 \mathscr Q:\mathscr D\longrightarrow\mathscr H^{16},\qquad
 \mathscr Q\psi=(\iota Q_\alpha\psi)_{\alpha=0}^{15}.
\end{equation}
Its source-normalized nonnegative quadratic form is
\begin{equation}\label{eq:core-form}
 q_0(\psi):=\frac1{16}\sum_{\alpha=0}^{15}
 \bigl\|\iota Q_\alpha\psi\bigr\|_{\mathscr H}^{2}
 =\frac1{16}\|\mathscr Q\psi\|^2_{\mathscr H^{16}}.
\end{equation}
\end{definition}

\begin{sourceaccount}[The source form is nonnegative]
For every $\psi\in\mathscr D$, $q_0(\psi)\geq0$.
\end{sourceaccount}
\begin{proof}
Each summand in \eqref{eq:core-form} is the square of a nonnegative
Hilbert norm. Their sum is nonnegative, and $1/16>0$. This elementary
fact is recorded as \texttt{AlgebraData.coreForm\_nonneg}; it is not a
spectral theorem.
\end{proof}
\formalrecord{\texttt{AlgebraData.bracket}, \texttt{structureConstant},
\texttt{gammaTwo}, \texttt{charge}, \texttt{chargeVector},
\texttt{coreForm}, and \texttt{coreForm\_nonneg} in
\texttt{Core.lean}. Atlas D-02, D-03.}

\section{Formal symmetry on test functions is not self-adjointness}
The source proves a core-symmetry theorem for each charge. A typical
integration-by-parts identity is
\begin{equation}\label{eq:core-symmetry}
 \langle Q_\alpha\psi,\phi\rangle_{\mathscr H}
 =\langle\psi,Q_\alpha\phi\rangle_{\mathscr H},
 \qquad\psi,\phi\in\mathscr D,
\end{equation}
with the source's inner-product convention. Compact support eliminates
the spatial boundary terms; self-adjointness of the original $\theta$
operators and the real gamma/color coefficients handle the finite
fermionic operations. But an operator may be symmetric and not
self-adjoint, because its adjoint may have a larger domain.

\begin{remark}[Two closures must not be confused]
The source defines an all-smooth-core charge graph closure in
$\mathscr H\times\mathscr H^{16}$. The \emph{physical} charge is instead
obtained by starting with the \emph{gauge-invariant smooth core}, then
closing the corresponding graph in the appropriate physical Hilbert
input. A result about the former does not automatically identify the
latter or its adjoint. This distinction became central in the historical
FCP physical-domain sequence, and Chapter~\ref{ch:hamiltonian} retains it.
\end{remark}

\section{What the October source-operator work actually established}
An important early task was to show that a parameterized source
\emph{deformed} charge at $(h,m)=(1,0)$ really is the original $Q_\alpha$,
not merely a formally similar family. The historical FCP G4--K9 chain
worked coefficient by coefficient, passing through the original matrix
commutators and finite index sums. Its accepted endpoint was stronger
than a numeric comparison: it proved equality of the source's
potential \emph{continuous linear maps}, then equality of the full
smooth-core charges and smooth-core forms.

\begin{theorem}[Original massless source equalities;
historically qualified]\label{thm:massless-identity}
With the original source color, gamma and fermionic data, the parameterized
source fields at $(h,m)=(1,0)$ satisfy, on the literal original smooth
core,
\begin{equation}\label{eq:massless-identification}
 Q^{\mathrm{def}}_\alpha(1,0)=Q_\alpha,
 \qquad q^{\mathrm{def}}_{1,0}=q_0.
\end{equation}
The equalities are those of the actual source definitions for every
admissible \texttt{AlgebraData}, not merely of a chosen finite sample.
\end{theorem}
\begin{proof}[Source-dependency explanation]
The original potential multiplier $Q^{\mathrm{def}}_\alpha(1,0)$ is
expanded in the exact ordered color/spin sums. Previously qualified
finite-sum identities (G4--K9A through K9C) establish the coefficient
regrouping. The G4--K9D source theorem identifies the \emph{full}
fermionic multiplier with the original bracket multiplier and therefore
identifies the differential charges on every test function. Applying
$\iota$, taking the original sixteen norm squares, and summing with the
same coefficient $1/16$ gives the form equality. The explanation records
the logic of the accepted formal chain; it does not silently treat a
finite floating-point check as a universal algebraic proof.
\end{proof}
\formalrecord{\path{FCP.BFSSSU2GaugeG4K9D.sourcePotentialMultipliers_equal},
\path{sourceUnconditionalCoreCharge_equal}, and
\path{sourceUnconditionalCoreForm_equal}; FCP strict
\href{https://github.com/etblink/Foundational-Convergence-Program/actions/runs/38011995220}{\#38011995220}
(9 October 2026 local). The integrated source core-energy interface is
\path{AlgebraData.massless_coreForm_energy}, strict
\href{https://github.com/etblink/math/actions/runs/42325087575}{\#42325087575}.
Atlas D-03, A-01.}

\section{Why the form will matter more than a formal expansion}
In later chapters we will seek an identity of the schematic form
\begin{equation}\label{eq:energy-schematic}
 q_0(\psi)=\tfrac12\|\nabla\psi\|^2
 +\int_{\mathcal X} V(x)\|\psi(x)\|^2\,dx
 +\int_{\mathcal X}\langle\psi(x),B(x)\psi(x)\rangle_{\mathcal F}\,dx,
\end{equation}
where $V$ is the original bosonic potential and $B$ is the \emph{unnormalized}
original fermionic field. The precise coefficients and color conventions
matter. Chapter~\ref{ch:energy} will distinguish source-accepted
identifications from finite diagnostics and conditional bounds. In
particular, \eqref{eq:energy-schematic} is not an assumption of a radial
spectral gap. The bosonic potential has flat directions, and a fermionic
field growing linearly in $\|x\|$ needs genuine control.

The next step is to impose the actual SU(2) gauge constraint and determine
the Hilbert space on which the physically relevant closed charge operator
acts. Only then can the adjoint-square construction be stated correctly.
